\documentclass[10pt, openany]{book}

\usepackage[T2A]{fontenc}    % Встановлює внутрішнє кодування шрифтів для кирилиці
\usepackage[utf8]{inputenc}  % Дозволяє розпізнавати символи UTF-8
\usepackage[ukrainian]{babel} % Налаштовує правила переносів та мовні команди
\usepackage{graphicx}
\usepackage[left=2cm,right=1.5cm,top=1.5cm,bottom=2cm]{geometry}
\usepackage{listings}
\usepackage{color}
\usepackage{setspace}
\usepackage{pifont}
\usepackage{amsmath}
\usepackage{indentfirst}

\definecolor{dkgreen}{rgb}{0,0.6,0}
\definecolor{gray}{rgb}{0.5,0.5,0.5}
\definecolor{mauve}{rgb}{0.58,0,0.82}

\lstset{frame=tb,
  language=Python,
  aboveskip=3mm,
  belowskip=3mm,
  showstringspaces=false,
  columns=flexible,
  basicstyle={\large\ttfamily},
  numbers=none,
  numberstyle=\tiny\color{gray},
  keywordstyle=\color{blue},
  commentstyle=\color{dkgreen},
  stringstyle=\color{mauve},
  breaklines=true,
  breakatwhitespace=true,
  tabsize=3
}

\definecolor{my-light-gray}{gray}{0.95}

% Створюємо окремий стиль для виводу терміналу
\lstnewenvironment{console}[1][]
{
    \lstset{
        backgroundcolor=\color{my-light-gray}, % Світло-сірий фон для консолі
        basicstyle=\ttfamily\small\color{black}, % Моноширинний шрифт
        frame=single,                    % Рамка навколо тексту
        frameround=tttt,                 % Закруглені кути рамки
        rulecolor=\color{my-light-gray},       % Колір рамки
        numbers=none,                    % Без нумерації рядків
        keepspaces=true,
        #1
    }
}
{}

\raggedbottom

\begin{document}
\pagestyle{plain}

\begin{large}
\thispagestyle{empty}

\begin{titlepage}
\begin{center}
        {\Large \textbf{Чернівецький Національний університет імені Юрія Федьковича \\
Факультет математики та інформатики}} \\[5cm]   % University at the top

        {\Large В. Косован} \\
        {\Large О. Матвій} \\
        {\Large Г. Мельник} \\
        {\Large В. Мельник} \\[5cm]                   % Names

        {\Huge \textbf{Інтелектуальний Аналіз Даних}} \\[0.5cm]

        {\Large Навчальний посібник} \\[9cm]                        % Title & type

        {\large Чернівці, 2026}                              % Bottom
    \end{center}
\end{titlepage}

% -------------------------------
% Second Page (info, disclaimer, etc.)
% -------------------------------
\newpage
\thispagestyle{empty}   % No page number here


УДК 004.738.5

\begin{center}
Рекомендовано до друку методичною радою факультету математики та інформатики \\
Чернівецького національного університету імені Юрія Федьковича \\
( протокол №  від       2026 року ) \\
\end{center}

Укладачі: \\
Василь Косован, кандидат фізико-математичних наук, асистент кафедри математичного моделювання. \\ 
Олександр Матвій, кандидат фізико-математичних наук, доцент кафедри математичного моделювання; \\
Галина Мельник, кандидат економічних наук, доцент кафедри прикладної математики та інформаційних технологій; \\
Василь Мельник, кандидат фізико-математичних наук, асистент кафедри математичного моделювання; \\


Інтелектуальний Аналіз Даних : навчальний посібник.. Укл.: Василь Косован, Олександр Матвій, Галина Мельник, Василь Мельник – 
Чернівці: Чернівецький національний університет ім. Юрія Федьковича, 2026. –  с.  \\

Посібник висвітлює базові принципи інтелектуального аналізу даних. Матеріали посібника підготовлені на 
основі методологічних досліджень авторів та відповідного курсу лекцій, містять основи теорії і практики аналізу даних.

Рецензенти: \\
Міночкін Дмитро Анатолійович, канд. щу сфері телекомунікацій, доцент кафедрителекомунікацій, 
Національний технічний університет України "Київський політехнічний інститут імені Ігоря Сікорського"\\
Семенов Віталій, сеніор інженер, SoftHouseGroup. \\

 
Для студентів, що здобувають освіту в галузі знань „Інформаційні технології” для спеціальностей 122 - 
“Комп’ютерні науки”, 124 - “Системний аналіз”, 113 - ”Прикладна математика” та інших спеціальностей, 
що пов'язані із вивченням комп'ютерних інформаційних технологій. УДК 004.738.5 \\


© Косован В.М., Матвій О.В., Мельник Г.В., Мельник В.С. \\

© Чернівецький національний університет, 2026 \\



    \vfill


\pagebreak
\tableofcontents

% \input ida-intro.tex

\input ida-1-data-collection.tex

\input ida-2-data-processing.tex

\input ida-3-data-visualization.tex

\input ida-4-lin-alg.tex

\input ida-5-nlp.tex

\input ida-6-prob.tex

\input ida-7-ml-intro.tex

\input ida-8-lin-classif.tex

\input ida-9-nonlinear-ml.tex

\input ida-10-ensemble-ml.tex

\input ida-11-unsupervised-ml.tex

\input ida-12-hypothesis-test.tex

\input ida-13-recommend-syst.tex

\input ida-14-ts.tex

\input ida-15-crysis-prognos.tex

\input ida-16-dl.tex


\end{large}
\end{document}